\section*{Zusammenfassung}
\begin{otherlanguage}{ngerman}
In-Context-Segmentierung sagt die Maske einer Zielstruktur anhand weniger
annotierter Beispiele vorher, sodass ein einzelnes Modell neue Klassen ohne
erneutes Training segmentieren kann. Bei medizinischen 3D-Bildern wird der
Rechenaufwand zur zentralen Einschränkung, da die Kosten dichter Attention
zwischen Ziel und Kontext quadratisch mit der Anzahl der Voxel wachsen. Diese
Arbeit entwickelt ein effizientes Modell zur 3D-In-Context-Segmentierung
entlang dreier Fragen: wie Bild- und Labelinformation kombiniert werden
sollte, wofür Rechenaufwand eingesetzt werden sollte und wie sich
Trainingsaufgaben über die annotierten Klassen hinaus gewinnen lassen.

Erstens schlagen wir Dual Attention vor, bei der Bild- und Label-Tokens in
getrennten Strömen verbleiben, die in jeder Schicht miteinander interagieren.
Bei gleichem Rechenaufwand und mit weniger Parametern übertrifft sie eine
einmalige Fusion beider Ströme in 12 von 13 Evaluationsgruppen, mit den
größten Zugewinnen bei neuen Zielstrukturen. Zweitens schlagen wir eine
Coarse-to-Fine-Kaskade vor, in der jede Stufe einen einzelnen Ausschnitt
verarbeitet, der anhand der Vorhersage der vorherigen Stufe platziert wird;
diese Vorhersage geht zudem als Prior in die Label-Tokens des Ziels ein. Das
Training mit den eigenen Vorhersagen des Modells als Prior übertrifft das
Training mit gestörten Ground-Truth-Masken. Drittens schlagen wir einen
Generator synthetischer Aufgaben vor, der prozedurale Formen in reale
Aufnahmen einzeichnet. Verglichen mit vollständig synthetischen Bildern als
Untergrund verbessert er die Ergebnisse auf drei von vier Läsionsdatensätzen.

Wir evaluieren das Modell auf TotalSegmentator und neun weiteren, nicht im
Training verwendeten Datensätzen und vergleichen es mit Medverse. Bei gleicher
Anzahl an Verfeinerungsschritten ist unser finales Modell 15-mal schneller und
im Mittel über die Datensätze ohne Läsionen genauer (Dice 0,615 gegenüber
0,552), während Medverse auf zwei dieser Datensätze genauer bleibt. Läsionen
bleiben für beide Modelle schwierig und sind das zentrale offene Problem.
\end{otherlanguage}
